% drafts/appendix-d-datatypes.tex — Algebraic Data Types
%
% Maps the implementation's type system to the mathematical structures.

The \textsc{FockMap} library implements the algebraic framework of this
paper using F\#'s type system.  This appendix documents the
correspondence between the mathematical objects and their
computational representations, and explains how the type system
enforces correctness properties statically.

%%====================================================================
\subsection*{Symbolic expression hierarchy}

The core algebraic types form a three-level hierarchy that mirrors the
structure of operator algebra:

\begin{center}
\begin{tabular}{@{}lll@{}}
\toprule
\textbf{Type} & \textbf{Mathematical object} & \textbf{Structure} \\
\midrule
\texttt{C<'u>} & Coefficient $\times$ unit & $c \cdot \hat{O}$ \\
\texttt{P<'u>} & Product of coefficients & $\prod_k c_k \hat{O}_k$ \\
\texttt{S<'u>} & Sum of products & $\sum_m P_m$ \\
\bottomrule
\end{tabular}
\end{center}

The type parameter \texttt{'u} (the ``unit'') is the operator alphabet.
By instantiating \texttt{'u} with different types, the same algebraic
machinery handles:

\begin{itemize}
  \item \texttt{C<LadderOperatorUnit>} --- fermionic $\ad{j}$, $\an{j}$
  \item \texttt{C<SectorLadderOperatorUnit>} --- sector-tagged mixed
    operators
  \item \texttt{C<Pauli>} --- single-qubit Pauli operators
\end{itemize}

This parameterization is not merely a software convenience.  It
ensures that the algebraic laws (associativity of products,
distributivity of sums over products, coefficient combination) are
implemented exactly once and shared across all operator types.
A bug in coefficient tracking would affect all instantiations
uniformly, making it detectable by any test suite.

%%====================================================================
\subsection*{Phase as a discriminated union}

The phase group $\Zi$ (\cref{def:phase-group}) is represented as a
four-constructor discriminated union rather than a complex number:

\begin{center}
\begin{tabular}{@{}ll@{}}
\toprule
\textbf{Constructor} & \textbf{Value} \\
\midrule
\texttt{P1} & $+1$ \\
\texttt{M1} & $-1$ \\
\texttt{Pi} & $+i$ \\
\texttt{Mi} & $-i$ \\
\bottomrule
\end{tabular}
\end{center}

Multiplication is a $4 \times 4$ pattern match with no arithmetic.
The conversion to a coefficient effect (\textsc{FoldIntoGlobalPhase})
uses only negation and the swap $(a + bi) \mapsto (-b + ai)$ ---
operations that are exact in IEEE~754.

This representation makes it \emph{impossible} for the implementation
to accumulate phase errors through repeated multiplication, because
there is no floating-point operation to round.  The exactness theorem
(\cref{thm:exactness}) is enforced by construction, not by careful
programming.

%%====================================================================
\subsection*{Pauli register and sequence types}

A single Pauli string is stored as a \texttt{PauliRegister}: a
coefficient (complex number) paired with an array of single-qubit
Pauli labels.  The \emph{signature} (the operator labels without the
coefficient, e.g., ``$XYZII$'') serves as a canonical string key.

A sum of Pauli strings is a \texttt{PauliRegisterSequence}: internally,
a \texttt{Dictionary<string, PauliRegister>} keyed by signature.  This
has two important consequences:

\begin{enumerate}
  \item \textbf{Like-term combination is automatic.}  Inserting a
    Pauli string whose signature already exists in the dictionary
    triggers coefficient addition.  If the sum is zero, the entry is
    removed.  No explicit ``simplify'' pass is needed.
  \item \textbf{Matching is exact.}  The dictionary key is a string,
    not a floating-point comparison.  Two Pauli strings combine if
    and only if they are operator-identical.
\end{enumerate}

%%====================================================================
\subsection*{The combining algebra interface}

The \texttt{ICombiningAlgebra<'op>} interface abstracts the
commutation relation used during normal ordering
(\cref{def:combining-algebra}).  It has a single method:

\begin{center}
\texttt{Combine}: $P \times C \to P[\;]$
\end{center}

\noindent
which takes an existing product and a new operator to append, and
returns an array of resulting products (one for simple reordering,
two when the fundamental relation fires).  The two implementations
differ only in the sign of the reordered term:

\begin{center}
\begin{tabular}{@{}lcc@{}}
\toprule
& \textbf{Fermionic} & \textbf{Bosonic} \\
\midrule
Same-index swap sign & $-1$ & $+1$ \\
Different-index swap sign & $-1$ & $+1$ \\
Identity term coefficient & $+1$ & $+1$ \\
\bottomrule
\end{tabular}
\end{center}

The normal-ordering algorithm is completely generic: it works by
iterating \texttt{Combine} and is unaware of which particle statistics
apply.  Selecting between CAR and CCR is a one-line type parameter
choice.

%%====================================================================
\subsection*{Sector tags for mixed systems}

Mixed fermion--boson operators carry a \texttt{ParticleSector} tag
(\texttt{Fermionic} or \texttt{Bosonic}) alongside their ladder
operator type.  The mixed normal-ordering pipeline
(\cref{sec:ext-mixed}) uses these tags to:

\begin{enumerate}
  \item Partition each product into fermionic and bosonic
    sub-sequences.
  \item Normal-order each sub-sequence with the appropriate algebra.
  \item Recombine via Cartesian product with multiplied coefficients.
\end{enumerate}

The type system prevents mixing: a fermionic operator cannot be
accidentally passed to the bosonic \texttt{Combine}, because their
\texttt{'op} type parameters are distinct.

%%====================================================================
\subsection*{Design rationale}

Three design choices deserve comment:

\paragraph{Why discriminated unions, not classes?}
Discriminated unions are \emph{closed}: adding a new constructor
requires modifying the type definition, which triggers compiler
errors at every incomplete pattern match.  This makes it impossible
to ``forget a case'' in the Pauli multiplication table or the phase
group law.  A class hierarchy, by contrast, allows silent extension
via subclassing, which is a source of encoding bugs.

\paragraph{Why generic types, not code generation?}
The \texttt{'u} type parameter means that new operator alphabets
(e.g., for anyonic statistics) can be supported by defining a new
type and a new \texttt{ICombiningAlgebra} instance --- no code
generator, no boilerplate.  The compiler verifies type consistency
at every call site.

\paragraph{Why dictionaries, not lists?}
Representing sums as signature-keyed dictionaries rather than
lists of terms ensures $O(1)$ amortized like-term lookup.  More
importantly, it eliminates a class of bugs where duplicate terms
persist through a computation and produce wrong Hamiltonian
coefficients.  The invariant ``at most one entry per signature''
is maintained by the data structure, not by a post-processing
simplification step.
